
The results for subtask A are shown in Tables \ref{tab:TaskATwitter} and \ref{tab:TaskASMS} for Twitter and for SMS messages, respectively;
those for subtask B are shown in Table \ref{tab:TaskBTwitter} for Twitter and in Table \ref{tab:TaskBSMS} for SMS messages.
%Moreover, Table \ref{tab:TaskATwitterFull} shows a detailed description of the characteristics of the systems that participated in subtask A, testing on Twitter data.
Systems are ranked by their scores for the constrained runs;
the ranking based on scores for unconstrained runs is shown as a subindex.

%The most common features were $n$-grams, sentiment lexicons, POS tags, capitalized words, URLs, and emoticons.
%Most systems performed normalization to clean up URLs and handle punctuation and @-mentions.
%Common external resources and tools include Weka, Stanford CoreNLP, MPQA, SentiWordNet, WordNet, NLTK, and sentiment lexicons.

For both subtasks, there were teams that only submitted results for the Twitter test set.
Some teams submitted both a constrained and an unconstrained version (e.g., AVAYA and teragram).
%and, overall, the unconstrained versions performed better.
As one would expect,
the results on the Twitter test set tended to be better than those on the SMS test set
since the SMS data was out-of-domain with respect to the training (Twitter) data.

Moreover, the results for subtask A were significantly better than those for subtask B, which shows that it is a much easier task,
probably because there is less ambiguity at the phrase-level.
%A summary of the results follows below.



\begin{table}[t!]
    \centering
    \begin{small}
    \begin{tabular}{l@{ }c@{ }c@{ }c@{ }c}
    \bf Run & \bf Const- & \bf Unconst- & \bf Use & \bf Super-\\
    & \bf rained & \bf rained & \bf Neut.? & \bf vised?\\
    \hline
    NRC-Canada & 69.02 &  & yes & yes\\
    GU-MLT-LT & 65.27 &  & yes & yes\\
    teragram & 64.86 & 64.86$_{(1)}$ & yes & yes\\
    BOUNCE & 63.53 &  & yes & yes\\
    KLUE & 63.06 &  & yes & yes\\
    AMI\&ERIC & 62.55 & 61.17$_{(3)}$ & yes & yes/semi\\
    FBM & 61.17 &  & yes & yes\\
    AVAYA & 60.84 & 64.06$_{(2)}$ & yes & yes/semi\\
    SAIL & 60.14 & 61.03$_{(4)}$ & yes & yes\\
    UT-DB & 59.87 &  & yes & yes\\
    FBK-irst & 59.76 &  & yes & yes\\
    nlp.cs.aueb.gr & 58.91 &  & yes & yes\\
    $^\diamond$UNITOR & 58.27 & 59.50$_{(5)}$ & yes & semi\\
    LVIC-LIMSI & 57.14 &  & yes & yes\\
    Umigon & 56.96 &  & yes & yes\\
    NILC\_USP & 56.31 &  & yes & yes\\
    DataMining & 55.52 &  & yes & semi\\
    $^\diamond$ECNUCS & 55.05 & 58.42$_{(6)}$ & yes & yes\\
    nlp.cs.aueb.gr & 54.73 &  & yes & yes\\
    ASVUniOfLeipzig & 54.56 &  & yes & yes\\
    SZTE-NLP & 54.33 & 53.10$_{(9)}$ & yes & yes\\
    CodeX & 53.89 &  & yes & yes\\
    Oasis & 53.84 &  & yes & yes\\
    NTNU & 53.23 & 50.71$_{(10)}$ & yes & yes\\
    UoM & 51.81 & 45.07$_{(15)}$ & yes & yes\\
    SSA-UO & 50.17 &  & yes & no\\
    SenselyticTeam & 50.10 &  & yes & yes\\
    UMCC\_DLSI\_(SA) & 49.27 & 48.99$_{(12)}$ & yes & yes\\
    bwbaugh & 48.83 & 54.37$_{(8)}$ & yes & yes/semi\\
    senti.ue-en & 47.24 & 47.85$_{(13)}$ & yes & yes\\
    SU-sentilab &  & 45.75$_{(14)}$ & yes & yes\\
    OPTWIMA & 45.40 & 54.51$_{(7)}$ & yes & yes\\
    REACTION & 45.01 &  & yes & yes\\
    uottawa & 42.51 &  & yes & yes\\
    IITB & 39.80 &  & yes & yes\\
    IIRG & 34.44 &  & yes & yes\\
    sinai & 16.28 & 49.26$_{(11)}$ & yes & yes\\
       \hline
    Majority Baseline & \multicolumn{2}{c}{29.19}  & N/A & N/A \\
    \hline
    \end{tabular}
    \end{small}
    \caption{\label{tab:TaskBTwitter}Results for subtask B on the Twitter dataset.
             The $^\diamond$ indicates a system submitted as constrained but which used additional Tweets or additional sentiment-annotated text
             to collect statistics that were then used as a feature.}
  \end{table}


\begin{table}[t!]
    \centering
    \begin{small}
    \begin{tabular}{l@{ }c@{ }c@{ }c@{ }c}
    \bf Run & \bf Const- & \bf Unconst- & \bf Use & \bf Super-\\
    & \bf rained & \bf rained & \bf Neut.? & \bf vised?\\
    \hline
    NRC-Canada & 68.46 &  & yes & yes\\
    GU-MLT-LT & 62.15 &  & yes & yes\\
    KLUE & 62.03 &  & yes & yes\\
    AVAYA & 60.00 & 59.47$_{(1)}$ & yes & yes/semi\\
    teragram &  & 59.10$_{(2)}$ & yes & yes\\
    NTNU & 57.97 & 54.55$_{(6)}$ & yes & yes\\
    CodeX & 56.70 &  & yes & yes\\
    FBK-irst & 54.87 &  & yes & yes\\
    AMI\&ERIC & 53.63 & 52.62$_{(7)}$ & yes & yes/semi\\
    $^\diamond$ECNUCS & 53.21 & 54.77$_{(5)}$ & yes & yes\\
    UT-DB & 52.46 &  & yes & yes\\
    SAIL & 51.84 & 51.98$_{(8)}$ & yes & yes\\
    $^\diamond$UNITOR & 51.22 & 48.88$_{(10)}$ & yes & semi\\
    SZTE-NLP & 51.08 & 55.46$_{(3)}$ & yes & yes\\
    SenselyticTeam & 51.07 &  & yes & yes\\
    NILC\_USP & 50.12 &  & yes & yes\\
    REACTION & 50.11 &  & yes & yes\\
    SU-sentilab &  & 49.57$_{(9)}$ & no & yes\\
    nlp.cs.aueb.gr & 49.41 & 55.28$_{(4)}$ & yes & yes\\
    LVIC-LIMSI & 49.17 &  & yes & yes\\
    FBM & 47.40 &  & yes & yes\\
    ASVUniOfLeipzig & 46.50 &  & yes & yes\\
    senti.ue-en & 44.65 & 46.72$_{(12)}$ & yes & yes\\
    SSA\_UO & 44.39 &  & yes & no\\
    UMCC\_DLSI\_(SA) & 43.39 & 40.67$_{(14)}$ & yes & yes\\
    UoM & 42.22 & 35.22$_{(15)}$ & yes & yes\\
    OPTWIMA & 40.98 & 47.15$_{(11)}$ & yes & yes\\
    uottawa & 40.51 &  & yes & yes\\
    bwbaugh & 39.73 & 43.43$_{(13)}$ & yes & yes/semi\\
    IIRG & 22.16 &  & yes & yes\\
       \hline
    Majority Baseline & \multicolumn{2}{c}{19.03}  & N/A & N/A \\
    \hline
    \end{tabular}
    \end{small}
    \caption{\label{tab:TaskBSMS}Results for subtask B on the SMS dataset.
             The $^\diamond$ indicates a system submitted as constrained but which used additional Tweets or additional sentiment-annotated text
             to collect statistics that were then used as a feature.}
  \end{table}


\subsection{Subtask A: Contextual Polarity}

Table \ref{tab:TaskATwitter} shows that subtask A, Twitter,
attracted 23 teams, who submitted 21 constrained and 7 unconstrained systems.
Five teams submitted both a constrained and an unconstrained system,
and two other teams submitted constrained systems that are on the boundary between being constrained and unconstrained.

One system was semi-supervised, and the rest were supervised.
The supervised systems used classifiers such as SVM (8 systems), Naive Bayes (7 systems), and Maximum Entropy (3 systems). Other approaches used include an ensemble of classifiers, manual rules, and a linear classifier. Two of the systems chose not to predict neutral as a possible classification label.

%while the rest tried to predict all three labels: positive, negative, and neutral.

The average F1-measure on the Twitter test set was 74.1\% for constrained systems and 60.5\% for unconstrained ones; this does not mean that using additional data does not help, it just shows that the best teams only participated with a constrained system.
NRC-Canada had the best constrained system with an F1-measure of 88.9\%, and AVAYA had the best unconstrained one with F1=87.4\%.

Table~\ref{tab:TaskASMS} shows the results for the SMS test set,
where 20 teams submitted 19 constrained and 7 unconstrained systems
(again, this included two teams that submitted boundary systems, marked accordingly).
The average F-measure on this test set was 70.8\% for constrained systems and 65.7\% for unconstrained systems. The best constrained system was that of GU-MLT-LT with an F-measure of 88.4\%, and AVAYA had the best unconstrained system with an F1 of 85.8\%.

\subsection{Subtask B: Message Polarity}

Table \ref{tab:TaskBTwitter} shows that subtask B, Twitter,
attracted 38 teams, who submitted 36 constrained and 15 unconstrained systems (and two boundary ones).

The average F1-measure was 53.7\% for the constrained and 54.6\% for the unconstrained systems.

These averages are much lower than those for subtask A,
which indicates that subtask B is harder,
probably because a message can contain parts expressing both positive and negative sentiment.

Once again, NRC-Canada had the best constrained system with an F1-measure of 69\%, followed by teragram, which had the best unconstrained system with an F1-measure of 64.9\%.

As Table \ref{tab:TaskBSMS} shows,
the average F1-measure on the SMS test set was 50.2\% for constrained and 50.3\% for unconstrained systems.
NRC-Canada had the best constrained system with an F1=68.5\%, and AVAYA had the best unconstrained one with F1-measure of 59.5\%.



\subsection{Overall}

Overall, the results achieved by the best teams were very strong, especially for the simpler subtask A:

\begin{itemize}
    \item F1=88.93, NRC-Canada on subtask A, Twitter;
    \item F1=88.37, GU-MLT-LT on subtask A, SMS;
    \item F1=69.02, NRC-Canada on subtask B, Twitter;
    \item F1=68.46, NRC-Canada on subtask B, SMS.
\end{itemize}

We can see that the strongest team overall was that of NRC-Canada, which was ranked first on three of the four conditions;
and it was second on subtask A, SMS.
There were two other teams that were strong across both tasks and on both test sets: GU-MLT-LT and AVAYA.
Three other teams, namely teragram, BOUNCE and KLUE, were ranked in the top-3 in at least one subtask and test set.
